\scanpage{6}
\mspage{\handwritten{(29}}
\begingroup
\makeatletter\tagsleft@true\makeatother
\[
\underline{\mathrm{Pt}}(X)^{\unsure{\circ|o|0}}\dashrightarrow
\underline{\mathrm{Pt}}(\handwritten{\widetilde}{X_{\text{ét}}})\ ,
\tag*{$(*)$}
\]
où le premier membre est défini dans 7.6.. Ceci posé, on a le

\noindent\uline{Théorème} 7.9.\quad Soit $X$ un préschéma. Le foncteur précédent $(*)$ est une \uline{an\deleted{\unsure{it|ti}}\handwritten{ti}-équivalence} de la catégorie des points géomètriques sur $X$ (avec comme morphismes les flèches des spécialisation) avec la catégorie des foncteurs fibres sur le topos étale $\handwritten{\widetilde}{X_{\text{ét}}}$.

Comme ce théorème ne servira plus dans la suite du séminaire nous nous bornons ici à une esquisse de démonstration, où nous nous permettrons certaines libertés avec les questions d'Univers.

a) \uline{Le foncteur envisagé est pleinement fidèle}. Pour tout $X\handwritten{\in}\operatorname{Ob} X_{\text{ét}}$, soit $X^{\handwritten{\unsure{\vee|v}}}: \handwritten{\widetilde}{X_{\text{ét}}} \dashrightarrow (\mathrm{Ens})$ défini par
\[
\handwritten{\check}{X}(F) = \operatorname{Hom}(X,F) = F(X).
\]

\noindent Notons que le foncteur fibre
\[
\handwritten{\varepsilon_{\xi}} : F \handwritten{\rightsquigarrow} F_{\handwritten{\xi}}
\]
peut s'écrire
\[
\handwritten{\varepsilon_{\xi}} \handwritten{\simeq} \mathop{\underset{\handwritten{\longrightarrow}}{\mathrm{lim}}}\limits_i X_i^{\handwritten{\vee}}\ ,
\]
où les $X_i$ sont des préschémas affines étales sur $X$ indexées par une certeine catégorie pseudo-filtrante (cf 4.3. et 4.5.). On a donc pour tout foncteur $\handwritten{\varphi}: \handwritten{\widetilde}{X_{\text{ét}}} \dashrightarrow (\mathrm{Ens})$:\margin{\tikz{\fill (0,0) circle (.018);\fill (.08,.12) circle (.018);}}
\[
\operatorname{Hom}(\handwritten{\varepsilon_{\xi}},\handwritten{\varphi}) \mathrel{\mathop{\text{-\kern.04em-}}\limits^{\handwritten{\sim}}} \mathop{\underset{\handwritten{\longleftarrow}}{\mathrm{lim}}}\limits_{\handwritten{i}} \operatorname{Hom}(X_i^{\handwritten{\vee}},\handwritten{\varphi})\ ,
\]
d'autre part il est bien connu que l'on a un isomorphisme bifonctoriel
\[
\operatorname{Hom}(X_i^{\handwritten{\vee}},\handwritten{\varphi}) \mathrel{\mathop{\text{-\kern.04em-}}\limits^{\handwritten{\sim}}} \handwritten{\varphi}(X_i)\ .
\]

\figure{\begin{tikzpicture}[x=1cm,y=1cm]
\draw (0,0) .. controls (.45,-.76) and (1.34,-1.79) .. (1.83,-2.47);
\end{tikzpicture}}
\endgroup

% archive sheet number: 5
